%% ============================================================
%% Chapter 7 Solutions
%% ============================================================

\section{Chapter 7: Algorithmic Differentiation}

%% ----------------------------------------------------------
\begin{solution}{7.1}{L3}

\solanswer

Annuity function: $A = P\,r(1+r)^n/[(1+r)^n - 1]$ at $P=1000$, $r=0.05$, $n=20$.

\solpart{a} \textbf{Numerical value.}
$(1.05)^{20} = 2.6533$. $A = 1000 \times 0.05(2.6533)/(2.6533-1)
= 1000 \times 0.13277/1.6533 = \mathbf{80.24}$ per period.

\solpart{b} \textbf{Analytic derivatives.}
Let $R = (1+r)^n$. Then $A = PrR/(R-1)$.
\[
  \frac{\partial A}{\partial P} = \frac{rR}{R-1} = A/P = 0.08024. \quad
  \SI{P}{A} = \frac{P}{A}\cdot\frac{A}{P} = 1.
\]
\[
  \frac{\partial A}{\partial r} = P\frac{(R-1)(R + rn R/r) - rR\cdot nR/r}{(R-1)^2}
  \text{ [chain rule; simplify to get~} \SI{r}{A}\text{]}.
\]
Using log-derivative for $\SI{r}{A}$: $\ln A = \ln P + \ln r + n\ln(1+r) - \ln[(1+r)^n - 1]$.
$\SI{r}{A} = \frac{d\ln A}{d\ln r} = 1 + \frac{nr}{1+r} - \frac{n r (1+r)^{n-1}}{(1+r)^n - 1}
\cdot r$.
At $r=0.05$, $n=20$: $\SI{r}{A} \approx 1 + 0.9524 - 0.9524\times2.6533/1.6533
\approx 1 + 0.9524 - 1.529 \approx \mathbf{0.424}$.

$\SI{n}{A} = \frac{n}{A}\cdot\frac{\partial A}{\partial n}
= r\ln(1+r)\cdot n\cdot\frac{(R-1) - R}{(R-1)^2/(R-1)}\cdots$
Numerically: $\SI{n}{A} \approx -\mathbf{0.468}$ (payments increase as $n$ decreases).

\solpart{c} \textbf{Forward mode for $\partial A/\partial r$ (seeding $\dot{r}=1$).}
$\dot{R} = n(1+r)^{n-1} = 20(1.05)^{19} = 50.539$.
$\dot{A} = P[\dot{r}R + r\dot{R})(R-1) - rR\dot{R}]/(R-1)^2$.
Numerically: matches analytic result $\partial A/\partial r \approx 680.2$.
$\SI{r}{A} = (r/A)\times680.2 = (0.05/80.24)\times680.2 \approx 0.424$. $\checkmark$

\solpart{d} \textbf{Reverse mode (all three in one pass).}
Seed $\bar{A} = 1$. Backpropagate: $\bar{R} = P[r(R-1)-rR]/(R-1)^2$,
then $\bar{r} = PR/(R-1) + \bar{R}\cdot n(1+r)^{n-1}$,
$\bar{P} = rR/(R-1)$, $\bar{n} = \bar{R}\cdot R\ln(1+r)$.
One backward pass yields all three. $\checkmark$

\solpart{e} \textbf{Operation counts.}
Forward mode (one parameter): $\sim$8 multiplications, 3 additions.
Reverse mode (all 3): $\sim$12 multiplications, 6 additions.
Forward-mode for all 3 parameters: $3\times8 = 24$ operations.
Reverse mode is $24/12 = 2\times$ cheaper for all parameters combined.

\solnote{Part (b): the analytic $\SI{r}{A}$ via the log-derivative is
the cleanest approach. Students who expand via the quotient rule often
make algebra errors. Encourage log-differentiation for compound expressions
involving products and ratios.}

\end{solution}

%% ----------------------------------------------------------
\begin{solution}{7.2}{L4}

\solanswer

\solpart{a} \textbf{Climate model, 800 parameters, 2 response functionals, 10 h/run.}
\textbf{Use adjoint AD.} With $m = 800$ and $r = 2$, FSE costs $801$ runs
($\approx 8010$ hours) while adjoint costs $3$ runs ($30$ hours).
If the model code supports AD (e.g., via Tapenade for Fortran),
implement the adjoint. If not, central-difference FSE is not feasible;
use adjoint or restrict to Morris screening ($\approx 8000$ evaluations
$= 80{,}000$ hours — still infeasible). This is a case where
developing the adjoint is the only practical option.

\solpart{b} \textbf{PK model, 5 parameters, 8 time points, analytic solutions.}
\textbf{Use analytic FSE (or log-derivatives directly).}
With analytic solutions, $\partial u/\partial p_j$ can be derived
symbolically. No AD or numerical FD is needed. Cost: trivial.
Using MATLAB \texttt{diff} or SymPy once gives all 5$\times$8 sensitivities.

\solpart{c} \textbf{Traffic simulation, 40 parameters, compiled C++, 3 min/run.}
\textbf{Use centered finite differences.}
AD is unavailable (compiled black box). Centered FD costs $81$ runs
$\times3$ min $= 243$ min $\approx 4$ hours. Feasible overnight.
Adjoint would require modifying the C++ source (possible but costly in developer time).
Morris screening ($15$ trajectories, $\approx 615$ runs) is also feasible
if nonlinearity is a concern.

\solpart{d} \textbf{MATLAB ODE45 model, 10 parameters, 1 \QOI, \texttt{if} statement.}
\textbf{Use centered finite differences.}
The \texttt{if} statement makes the model non-differentiable at the quarantine
threshold; forward-mode AD and the adjoint will both fail at parameter values
that trigger the discontinuity. Centered differences detect the discontinuity
as a large error indicator (the left and right estimates diverge near threshold).
Cost: $21$ runs. Use the U-curve diagnostic to identify non-smooth parameters.

\solnote{Part (d) is the key discriminating item: students who recommend AD
or adjoint without noting the \texttt{if}-statement issue have not read the
question carefully. Chapter 8 covers this exact failure mode. The correct
approach for non-smooth models is FD (which still works, though less accurately
near the threshold) with careful step-size checking.}

\end{solution}

%% ----------------------------------------------------------
\begin{solution}{7.3}{L4}

\solanswer

\solpart{a} \textbf{dlarray for $R_0(k,\beta,\tau) = k\beta\tau$.}
\begin{verbatim}
k = dlarray(5); beta = dlarray(0.06); tau = dlarray(7);
R0 = k * beta * tau;
[dR0_dk, dR0_dbeta, dR0_dtau] = dlgradient(R0, k, beta, tau);
% dR0_dk = beta*tau = 0.42; dR0_dbeta = k*tau = 35; dR0_dtau = k*beta = 0.3
\end{verbatim}

\solpart{b} \textbf{Verification.}
$\partial R_0/\partial k = \beta\tau = 0.42$ $\checkmark$;
$\partial R_0/\partial\beta = k\tau = 35$ $\checkmark$;
$\partial R_0/\partial\tau = k\beta = 0.3$ $\checkmark$.
Normalized: $\SI{k}{R_0} = (k/R_0)\cdot0.42 = (5/2.1)\cdot0.42 = 1$ $\checkmark$.

\solpart{c} \textbf{SIR endemic equilibrium.}
At the endemic equilibrium $\bar{I} > 0$: numerically solve
$\vec{f}(\vec{u};\beta) = \vec{0}$ for $\bar{I}(\beta)$.
Since \texttt{ode45} is not natively differentiable via \texttt{dlarray},
implement a Newton iteration inside a \texttt{dlarray} context:
\begin{verbatim}
beta_dl = dlarray(0.06);
% Newton iteration (differentiable):
u = dlarray([999; 1; 0]);  % initial guess
for k = 1:20
    F = sir_rhs_dl(u, beta_dl);  % model RHS as dlarray function
    J = sir_jac_dl(u, beta_dl);
    u = u - J \ F;
end
dIbar_dbeta = dlgradient(u(2), beta_dl);
\end{verbatim}
Compare with FSE result from Chapter 4 ($\SI{\beta}{\bar{I}}$ at equilibrium).

\solpart{d} \textbf{Applying dlgradient through ode45.}
\texttt{ode45} is not a differentiable MATLAB function: it uses
internal adaptive stepping, logical branching, and memory allocation
that \texttt{dlarray} cannot trace through the computation graph.
Calling \texttt{dlgradient} after \texttt{ode45} will produce an error
or silently return zero gradients.
\textbf{Remedy:} Either (1) implement the ODE solve as a differentiable
fixed-step method (e.g., explicit Euler or RK4 loop in \texttt{dlarray}),
or (2) use the adjoint ODE approach (Chapter 6), which bypasses the
need to differentiate through the ODE solver itself.

\solnote{Part (d) is the central limitation of black-box AD applied to
ODE solvers. The forward-mode approach (seeding $\dot{\beta}=1$ through
the ODE RHS) works only for fixed-step explicit integrators. The adjoint
approach of Chapter 6 circumvents this entirely.}

\end{solution}

%% ----------------------------------------------------------
\begin{solution}{7.4}{L5}

\solanswer

\textbf{Dengue model SA design: 8 state variables, 14 parameters, 2 QOIs, 2 min/run.}

\solpart{a} \textbf{Method: analytic adjoint ODE.}
With $m = 14$ and $r = 2$, FSE costs $15$ runs ($30$ min) while adjoint
costs $3$ runs ($6$ min). For a published dengue model with known equations,
deriving the $8\times8$ Jacobian analytically is feasible ($\sim1$ day of
derivation). Adjoint is preferred for efficiency and for providing exact
sensitivities without step-size issues. If derivation time is a constraint,
centered FD at $29$ runs ($58$ min) is also acceptable.

\solpart{b} \textbf{Computation time.} Adjoint: 1 forward + 1 adjoint
(for $r = 2$, need 2 adjoint solves) = 3 runs $\times$ 2 min = \textbf{6 minutes.}
Centered FD: $29$ runs $\times$ 2 min = \textbf{58 minutes.}

\solpart{c} \textbf{Validation.}
Verify $\SI{k}{R_0} = \SI{\beta_h}{R_0} = 1$ analytically against adjoint output.
Compare adjoint sensitivities with centered FD for all 14 parameters;
require agreement to 4 significant figures. Test the adjoint code on a
single-compartment analytic submodel with known sensitivities.

\solpart{d} \textbf{Handling structural correlation between $\sigma_v$ and $\sigma_h$.}
Report standard (marginal) SIs and total-derivative SIs. The coupling
coefficient is estimated from the epidemiological literature or from
field measurements of the joint distribution of biting rates. If
$\sigma_v$ increases by $x\%$, $\sigma_h$ increases by $c_{vh}\cdot x\%$.
The total-derivative SI of each QOI with respect to $\sigma_v$ accounts
for this induced change in $\sigma_h$. Report both sets side-by-side and
flag the structural correlation explicitly in the methods section.

\solpart{e} \textbf{Presentation for public health audience.}
Two tornado plots (one per QOI), with parameters labeled by plain-language
intervention descriptions: ``Reducing average mosquito biting rate by 10\%
reduces disease prevalence by approximately X\%.'' Use a 3-color coding:
green (large positive SI — increasing this parameter helps reduce burden),
red (large negative SI — decreasing this parameter reduces burden),
gray (small SI — negligible influence). Limit to top 5 parameters
per QOI for clarity.

\solnote{This is a realistic applied SA design exercise. Full credit requires
all five components with specificity. The structural correlation handling (d)
is the most commonly missed; students who only say ``account for correlation''
without specifying how receive partial credit.}

\end{solution}

%% ----------------------------------------------------------
\begin{solution}{7.5}{L6}

\solanswer

\textbf{The most likely technical cause:} Incorrect treatment of the
adaptive step-size logic during reverse-mode AD of the ODE solver.

\solpart{a} \textbf{Root cause.}
Adaptive-step-size ODE solvers (like \texttt{ode45}) use internal
comparisons and branching to decide step sizes. When reverse-mode AD
is applied through the ODE solver, the branching decisions made during
the forward pass are replayed during the backward pass. For parameters
that appear in the mortality rate expressions (which affect stability
and thus step-size selection), a perturbation can change which branches
are taken, making the reverse-mode computation inconsistent with the
actual forward computation. This manifests as 10--20\% errors that
are largest for parameters affecting stability (mortality rates).

\solpart{b} \textbf{Two diagnostic tests.}

\textit{Test 1: Fixed-step integrator comparison.}
Re-implement the ODE solve with a fixed-step explicit RK4 method.
Apply reverse-mode AD to this fixed-step implementation. If the
discrepancy disappears, the adaptive step-size branching is the cause.

\textit{Test 2: Finite-difference consistency.}
For the discrepant parameters, compute FD sensitivities at $h = 10^{-4}$
and $h = 10^{-5}$. If these agree with each other but not with the
AD result, the AD implementation is wrong; if they also disagree,
the model may be genuinely discontinuous near the nominal parameter values.

\solpart{c} \textbf{Two remedies and trade-offs.}

\textit{Remedy 1: Fixed-step AD-compatible integrator.}
Replace the adaptive solver with a fixed-step RK4 inside the AD tape.
\textit{Trade-off:} Fixed-step integrators require choosing a step size
that ensures accuracy for all parameter values; for stiff models,
this may require very small steps and high computational cost.
Advantage: mathematically clean, gradients are exact for the fixed-step
approximation.

\textit{Remedy 2: Discrete adjoint (checkpointing).}
Use the adjoint ODE approach (Chapter 6) with the adaptive forward solver
stored at checkpoints. The adjoint ODE is integrated backward using
the stored checkpoints, avoiding the need to differentiate through
the ODE solver's branching logic.
\textit{Trade-off:} Requires implementing the adjoint ODE separately
(additional derivation and coding); storage grows with the number of checkpoints.
Advantage: handles adaptive steps naturally and is more memory-efficient.

\solpart{d} \textbf{When to conclude genuine SA result.}

If the FD sensitivities computed at two consistent step sizes ($h$ and $h/2$)
agree with each other within 1\% but differ from the AD result by 10--20\%,
the AD implementation is suspect. However, if FD estimates also show
high variability (differing by $>5\%$ as $h$ is halved), this suggests
the model is genuinely non-smooth and the large discrepancy reflects
real behavior: the model is in a sensitive regime where small parameter
changes cause qualitative changes in dynamics. In this case, the adjoint
and FD are both locally valid; the discrepancy is genuine and should
be reported as a finding (the model is near a threshold for these parameters).

\solnote{This exercise mirrors real-world challenges in gradient-based
optimization of climate models. The root cause (adaptive branching in
reverse-mode AD) is well-documented in the scientific computing literature.
Students who identify only ``there is a bug'' without specifying the mechanism
receive partial credit. The key insight is that AD through adaptive-step
ODE solvers is non-trivial, and the discrete adjoint (Chapter 6's approach)
is the standard remedy.}

\end{solution}

%% ============================================================
%% Chapter 8 Solutions
%% ============================================================

\section{Chapter 8: Caveat Emptor}

%% ----------------------------------------------------------
\begin{solution}{8.1}{L3}

\solanswer

$g(u;p) = u^3 - 3pu + 2 = 0$ (same as Ex 4.2, now with bifurcation focus).

\solpart{a} \textbf{Real roots at $p = 1$ and most rapidly changing.}
Roots: $u_1 = 1$ (double), $u_2 = -2$. The double root $u_1 = 1$
changes most rapidly — its sensitivity is undefined (infinite), as shown
in Ex 4.2. It is at a fold bifurcation point.

\solpart{b} \textbf{$\SI{p}{u^*}$ at each root.}
As computed in Ex 4.2: $\SI{p}{u_1}$ is undefined (singular FSE);
$\SI{p}{u_2} = 1/3$.

\solpart{c} \textbf{Root curve and bifurcation.}
The real roots satisfy $u^3 - 3pu + 2 = 0$. By Cardano's formula or
numerical continuation, the two positive roots merge at the fold bifurcation
where $\partial g/\partial u = 3u^2 - 3p = 0$, i.e., $u = \sqrt{p}$.
Substituting: $p^{3/2} - 3p\sqrt{p} + 2 = 0 \Rightarrow -2p^{3/2} + 2 = 0
\Rightarrow p^* = 1$. The two roots $u_1 = 1$ and $u_2 = -2$ merge
at $p = 1$ (actually only the double root merges; $u_2 = -2$ persists).
Careful: the cubic has the double root at $u = 1$ for $p = 1$.
For $p < 1$, the double root splits into two complex roots (no real positive root
near 1), and for $p > 1$, two distinct real roots appear near $u = 1$.
The bifurcation is at $\mathbf{p = 1}$.

\solpart{d} \textbf{FSE prediction at bifurcation.}
At $p = 1$, $\partial g/\partial u = 0$: the FSE $(\partial g/\partial u)\,(\partial u/\partial p) = -\partial g/\partial p$ has a zero coefficient, so $\partial u/\partial p = \infty$.
Finite-difference verification: $u(p=1.001)$ vs $u(p=0.999)$ shows very
large differences near the double root, confirming infinite sensitivity.
The FD estimate diverges as $p \to 1$, consistent with a fold bifurcation.

\solnote{This exercise directly connects Chapter 4 (implicit FSE) to
Chapter 8 (bifurcations). The key point: the FSE singularity \emph{is}
the bifurcation. Students who answer only ``sensitivity is large'' without
connecting to the eigenvalue of the Jacobian receive partial credit.}

\end{solution}

%% ----------------------------------------------------------
\begin{solution}{8.2}{L3}

\solanswer

\solpart{a} \textbf{$\SI{\beta}{I(T)}$ at four values, $T = 90$ days.}

Computed via centered FD with $h = 6\times10^{-6}\beta^*$:

\begin{center}
\begin{tabular}{lcccc}
\toprule
$\beta$ & $R_0 = k\beta\tau$ & $I(90)$ & $\SI{\beta}{I(90)}$ (approx.) \\
\midrule
0.04 & 1.40 & $\sim150$ & $\sim+3.2$ \\
0.06 & 2.10 & $\sim480$ & $\sim+1.8$ \\
0.08 & 2.80 & $\sim620$ & $\sim+1.2$ \\
0.12 & 4.20 & $\sim720$ & $\sim+0.7$ \\
\bottomrule
\end{tabular}
\end{center}

The SI \emph{decreases} as $R_0$ increases past 1: at higher $R_0$,
almost everyone has been infected by day 90, and $I(90)$ is in the
declining tail, less sensitive to $\beta$.

\solpart{b} \textbf{$\beta = 0.02$ ($R_0 = 0.7 < 1$).}
No epidemic occurs: $I(90) \approx I(0) = 1$ (infection dies out
exponentially). $\SI{\beta}{I(90)}$ is small and positive near zero.
\textbf{Qualitative change:} the SI is negligibly small for $R_0 < 1$
because no epidemic amplification occurs, but large and positive
for $R_0 > 1$. Near $R_0 = 1$ ($\beta \approx 0.029$), the SI
would be extremely large — the threshold crossing amplifies any
perturbation enormously.

\solpart{c} \textbf{Bifurcation diagram.}
The endemic equilibrium $\bar{I}^*$ bifurcates from $\bar{I}^* = 0$
at $R_0 = 1$ (transcritical bifurcation). For $R_0 < 1$: $\bar{I}^* = 0$
(disease-free equilibrium is stable). For $R_0 > 1$: $\bar{I}^* > 0$
(endemic equilibrium exists and is stable). The bifurcation diagram shows
$\bar{I}^*$ vs.\ $\beta$ as a forward (supercritical) transcritical bifurcation.

\solnote{The key finding from part (a) is that $\SI{\beta}{I(T)}$ \emph{decreases}
as $\beta$ increases — counterintuitive if one expected larger $\beta$ to
mean larger sensitivity. The reason: at high $\beta$, the epidemic has
already run its course by day 90, and $I(90) \approx 0$ regardless of
small perturbations to $\beta$. This illustrates why SA at a single time
point can be misleading (Chapter 4, Ex 4.8).}

\end{solution}

%% ----------------------------------------------------------
\begin{solution}{8.3}{L4}

\solanswer

\solpart{a} \textbf{Condition number and error vs.\ $\epsilon$.}

MATLAB code:
\begin{verbatim}
u_exact = [1;1;1];
eps_vals = [1, 0.1, 0.01, 0.001, 0.0001];
for k = 1:length(eps_vals)
    ep = eps_vals(k);
    A = [1 2 3; 4 5 6; 7 8 9+ep];
    b = A * u_exact;
    kappa(k) = cond(A);
    err(k) = norm(A\b - u_exact);
end
\end{verbatim}

Expected results (approximate):

\begin{center}
\begin{tabular}{rcc}
\toprule
$\epsilon$ & $\kappa(A)$ & $\|\hat{\vec{u}} - \vec{u}\|$ \\
\midrule
1.0    & $\sim 30$    & $\sim 10^{-14}$ \\
0.1    & $\sim 300$   & $\sim 10^{-13}$ \\
0.01   & $\sim 3000$  & $\sim 10^{-12}$ \\
0.001  & $\sim 3\times10^4$ & $\sim 10^{-11}$ \\
0.0001 & $\sim 3\times10^5$ & $\sim 10^{-10}$ \\
\bottomrule
\end{tabular}
\end{center}

\solpart{b} \textbf{Log-log slope.}
The relationship is $\|\hat{\vec{u}} - \vec{u}\| \propto \kappa(A)\,\varepsilon_{\text{mach}}$.
On log-log axes, the slope is $+1$ (error grows linearly with condition number).

\solpart{c} \textbf{Implication for large SI values.}

A large SI for a nearly singular model reflects the ill-conditioning of the
sensitivity computation, not necessarily a genuine large sensitivity of the
physical model. If the Jacobian $J_F$ of the model equations is nearly
singular (large $\kappa(J_F)$), the FSE $J_F\,(\partial\vec{u}/\partial p_j) = \vec{f}_j$
will amplify any numerical errors in $\vec{f}_j$ by a factor of $\kappa(J_F)$.
The reported SI may be $\kappa(J_F)$ times larger than the true value.
\textit{Practical implication:} always report $\kappa(J_F)$ alongside
large SI values. An SI of $50$ with $\kappa(J_F) = 100$ should be
flagged as potentially unreliable; an SI of $50$ with $\kappa(J_F) = 2$ is trustworthy.

\solnote{The slope of 1 in part (b) is the classical condition-number
error bound: $\|\delta\vec{u}\|/\|\vec{u}\| \leq \kappa(A)\|\delta\vec{b}\|/\|\vec{b}\|$.
Students who compute the slope empirically and report it as $\approx 1$
without citing the bound receive full credit; students who derive the bound
analytically receive bonus credit.}

\end{solution}

%% ----------------------------------------------------------
\begin{solution}{8.4}{L5}

\solanswer

\textbf{Advisory report: SA near the epidemic threshold ($R_0 = 1.04$).}

\solpart{a} \textbf{What local SA can and cannot tell you.}

Local SA provides the derivative of each \QOI with respect to each parameter
at the current best-fit values. With $R_0 = 1.04$, the model is very
close to the epidemic threshold. Local SA \emph{can} tell you which
parameters would, if perturbed slightly, have the largest marginal effect
on disease burden. What it \emph{cannot} tell you: whether a parameter
change of realistic magnitude (say 10--20\%) would push $R_0$ below 1
(eliminating the epidemic entirely) or keep it above 1. Near the threshold,
the model's behavior changes qualitatively, and local SA captures only the
linearized behavior in an infinitesimally small neighborhood.

\solpart{b} \textbf{Specific numerical problems near $R_0 = 1$.}

Near the transcritical bifurcation at $R_0 = 1$, the Jacobian of the
model at the endemic equilibrium has an eigenvalue near zero. This causes:
(1) the FSE linear system $J_F\,(\partial\vec{u}/\partial p) = \vec{f}$
to be nearly singular, amplifying roundoff errors by $\kappa(J_F)$;
(2) time-dependent SIs to grow without bound near the threshold crossing;
(3) the adjoint ODE to have a slowly decaying mode (near-zero eigenvalue),
making backward integration unstable over long time horizons.
Practitioners should report $\kappa(J_F)$ and check that SI values are
stable to perturbations of the nominal parameter set within the confidence region.

\solpart{c} \textbf{Additional analyses recommended.}

\begin{enumerate}[nosep]
\item \textit{Bifurcation diagram:} Plot the endemic equilibrium $I^*$ vs.\
  $R_0$ (or vs.\ the most controllable parameter). Determine how large a
  reduction in $\beta$ or $k$ is needed to push $R_0$ below 1.
  This gives a specific, actionable target rather than a marginal SI value.
\item \textit{SA at multiple $R_0$ values:} Compute SI at $R_0 = 0.9$,
  $1.0$, $1.04$, $1.2$. Report how the rankings change. If the ranking
  is stable across this range, the local SA at $R_0 = 1.04$ is informative.
  If not, flag the instability.
\item \textit{Finite-change analysis:} Compute the actual change in peak $I$
  for a $10\%$ reduction in each parameter. Report these directly alongside
  the SI values. For parameters near the threshold, the finite change may
  be qualitatively different from the SI prediction.
\end{enumerate}

\solpart{d} \textbf{Explaining bifurcation to a policy maker.}

Think of the epidemic threshold as a tipping point. Below it ($R_0 < 1$),
each infected person infects on average less than one other person, and
the outbreak fades on its own. Above it ($R_0 > 1$), each infected person
infects more than one other person, and the outbreak grows into an epidemic.
At $R_0 = 1.04$, we are just above the tipping point. Even a small
reduction in contact rates or transmission probability — say $5\%$ —
could push us below the tipping point, causing the epidemic to die out
rather than sustain itself. This is the most important insight from the
sensitivity analysis: at $R_0$ close to 1, small interventions can have
disproportionately large effects. But it also means our predictions are
highly uncertain: if the true $R_0$ is $0.96$ rather than $1.04$, the
epidemic might already be declining naturally.

\solnote{This is a realistic scenario that every applied mathematical modeler
faces. The report format requirement is deliberate: public health advisors
need concise, clear guidance, not a technical treatise. Deduct credit for
responses that are technically correct but not accessible to a non-mathematical
reader. The bifurcation explanation (part d) should use no equations.}

\end{solution}

%% ----------------------------------------------------------
\begin{solution}{8.5}{L6}

\solanswer

\textbf{Three methodological concerns with the malaria SA ($R_0 = 1.15$, all SI $< 0.1$):}

\textbf{Concern 1: Small SIs near $R_0 = 1$ may be numerical artifacts.}

\textit{Problem:} With $R_0 = 1.15$, the model is moderately above the
threshold, but close enough that the endemic equilibrium Jacobian may be
ill-conditioned. Small SI values could result from condition-number amplification
rather than genuine insensitivity.

\textit{Evidence needed:} Report $\kappa(J_F)$ at the endemic equilibrium
and verify that $|\text{SI}| \gg \kappa(J_F)\,\varepsilon_{\text{mach}}$
for each parameter. If any SI is within a factor of $\kappa(J_F)$ of machine
epsilon, that SI is numerically unreliable.

\textit{Alternative analysis:} Compute SIs of $R_0$ itself (analytically)
and compare sign and magnitude with the equilibrium SIs. If $\SI{p_j}{R_0}$
is large but $\SI{p_j}{\bar{I}}$ is small, the discrepancy suggests
numerical issues at the equilibrium.

\textbf{Concern 2: SA of the endemic equilibrium does not address
transient dynamics or the probability of outbreak onset.}

\textit{Problem:} If malaria is currently at or near the endemic equilibrium,
the relevant policy question is ``which intervention would most reduce
endemic prevalence?'' But if the model is used for outbreak forecasting
(a different but common use), the relevant QOI is the peak or cumulative
burden during a transient following an introduction event. SA of the
endemic equilibrium cannot answer the latter question.

\textit{Evidence needed:} Clarify the model's purpose. If forecasting
an outbreak, compute SA for the transient QOI (e.g., peak $I$ at $T = 180$
days) in addition to the endemic equilibrium.

\textit{Alternative analysis:} Compute time-dependent SIs $\SI{p_j}{I}(t)$
and report whether the parameter ranking changes between the transient
and endemic phases.

\textbf{Concern 3: ``All indices less than 0.1'' does not imply robustness
unless the uncertainty ranges are accounted for.}

\textit{Problem:} A local SI of 0.1 means a 1\% change in the parameter
produces a 0.1\% change in prevalence. But if a parameter is uncertain
by $\pm50\%$ (plausible for entomological parameters in malaria models),
the expected change in prevalence from this parameter alone is $\approx 5\%$.
Reporting only SI values without the uncertainty range gives an incomplete
picture.

\textit{Evidence needed:} Report sigma-normalized SIs or one-way sensitivity
ranges: the change in $\bar{I}$ when each parameter varies over its full
uncertainty range. Parameters with small SI but large uncertainty range
may still dominate the output variance.

\textit{Alternative analysis:} Compute Sobol first-order and total indices
across the plausible parameter uncertainty ranges. A parameter with local
SI = 0.1 and uncertainty range $\pm80\%$ may have Sobol $S_1 = 0.3$.

\solnote{Full credit requires all three concerns with (a), (b), (c) for each.
The three concerns should be genuinely distinct: numerical reliability
(Concern 1), scope of the SA (Concern 2), and accounting for uncertainty
ranges (Concern 3). Students who only note ``local SA doesn't capture
nonlinearity'' receive partial credit for Concern 3 without the sigma-normalized
index framing. The claim ``all indices less than 0.1'' is particularly dangerous
near a threshold --- this is the chapter's core message.}

\end{solution}

%% ============================================================
%% Chapter 9 Solutions
%% ============================================================

\section{Chapter 9: Global Sensitivity Analysis}

%% ----------------------------------------------------------
\begin{solution}{9.1}{L3}

\solanswer

Given: $\SI{\tau_E}{R_0} = 0.6$, $\SI{\tau_I}{R_0} = 1.2$,
$\SI{\beta}{R_0} = 1.0$, $\SI{k}{R_0} = 1.0$, nominal $R_0^* = 2.0$.
Uncertainties: $\sigma_{\tau_E} = 2$ d, $\sigma_{\tau_I} = 3$ d,
$\sigma_\beta = 0.01$, $\sigma_k = 0.8$.
Nominal values: $\tau_E^* = 5.5$ d, $\tau_I^* = 7$ d, $\beta^* = 0.06$,
$k^* = 5$.

\solpart{a} \textbf{Output standard deviation (delta method).}

$\sigma_{R_0}^2 = \sum_j \left(\frac{\partial R_0}{\partial p_j}\right)^2 \sigma_{p_j}^2
= \sum_j \left(\frac{R_0^*}{p_j^*}\,\SI{p_j}{R_0}\right)^2 \sigma_{p_j}^2$.

\begin{align*}
\left(\frac{R_0^*}{\tau_E^*}\,\SI{\tau_E}{R_0}\right)^2\sigma_{\tau_E}^2
&= \left(\frac{2}{5.5}\times0.6\right)^2\times4 = (0.218)^2\times4 = 0.190\\
\left(\frac{R_0^*}{\tau_I^*}\,\SI{\tau_I}{R_0}\right)^2\sigma_{\tau_I}^2
&= \left(\frac{2}{7}\times1.2\right)^2\times9 = (0.343)^2\times9 = 1.059\\
\left(\frac{R_0^*}{\beta^*}\,\SI{\beta}{R_0}\right)^2\sigma_\beta^2
&= \left(\frac{2}{0.06}\times1.0\right)^2\times(0.01)^2 = (33.3)^2\times10^{-4} = 0.111\\
\left(\frac{R_0^*}{k^*}\,\SI{k}{R_0}\right)^2\sigma_k^2
&= \left(\frac{2}{5}\times1.0\right)^2\times(0.8)^2 = (0.4)^2\times0.64 = 0.102
\end{align*}
$\sigma_{R_0}^2 = 0.190 + 1.059 + 0.111 + 0.102 = 1.462$,
$\sigma_{R_0} \approx \mathbf{1.21}$.

\solpart{b} \textbf{Sigma-normalized indices.}

$\tilde{S}_{p_j} = (R_0^*/p_j^*)\,|\SI{p_j}{R_0}|\,\sigma_{p_j}/\sigma_{R_0}$:

\begin{center}
\begin{tabular}{lcc}
\toprule
Parameter & $\tilde{S}$ (sigma-norm.) & Contribution (\%) \\
\midrule
$\tau_I$ & $(2/7)(1.2)(3)/1.21 = 0.851$ & $58\%$ \\
$\tau_E$ & $(2/5.5)(0.6)(2)/1.21 = 0.360$ & $13\%$ \\
$\beta$  & $(2/0.06)(1.0)(0.01)/1.21 = 0.276$ & $8\%$ \\
$k$      & $(2/5)(1.0)(0.8)/1.21 = 0.264$ & $7\%$ \\
\bottomrule
\end{tabular}
\end{center}

\solpart{c} \textbf{Tornado plot.} Sorted by $\tilde{S}$: $\tau_I > \tau_E > \beta \approx k$.

\solpart{d} \textbf{Ranking comparison.}
Local SI ranking: $\tau_I$ and $k$ tied at 1.2 and 1.0; $\beta$ at 1.0; $\tau_E$ at 0.6.
Sigma-normalized ranking: $\tau_I \gg \tau_E > \beta \approx k$.

$\tau_I$ remains dominant in both, but its lead is much larger when
uncertainty is accounted for ($\sigma_{\tau_I} = 3$ d is large).
$k$ drops from second to last: despite SI $= 1.0$, its uncertainty
$\sigma_k = 0.8$ relative to $k^* = 5$ gives a coefficient of variation
of only $16\%$, while $\tau_E$ has $\sigma_{\tau_E}/\tau_E^* = 36\%$.
\textbf{$k$ changed rank most dramatically} because it has a large SI
but small relative uncertainty.

\solnote{The sigma-normalized index is the key concept here: it combines
the sensitivity (SI) with the actual uncertainty range to produce a
practical priority ranking. Students who compute only SIs without
accounting for uncertainty miss the point of this exercise.}

\end{solution}

%% ----------------------------------------------------------
\begin{solution}{9.2}{L3}

\solanswer

\solpart{a} \textbf{LHS implementation and comparison.}

\begin{verbatim}
rng(42);
N = 100; m = 3;
X_lhs = lhs_sample(N, m, zeros(1,m), ones(1,m));
X_rnd = rand(N, m);
\end{verbatim}

\solpart{b} \textbf{Scatter plots.} The LHS scatter plots show more
uniform coverage with fewer clusters; the random sample shows visible gaps
and clusters in all three pairwise projections.

\solpart{c} \textbf{Mean and standard deviation.}
Both LHS and random sampling should give means near $0.5$ and
standard deviations near $1/\sqrt{12} \approx 0.289$ for a uniform
distribution on $[0,1]$. LHS is typically closer to both exact values
because the stratification ensures coverage of the full range. For $N = 100$,
the improvement in mean accuracy is $\sim$30--50\% and in standard deviation
$\sim$20--40\%.

\solpart{d} \textbf{Histogram comparison.}
The LHS histogram shows a much flatter, more uniform distribution
(each of the 100 strata has exactly one sample). The random sample
histogram shows the typical Poisson variation (some bins have 0 or 3+ samples).
For sensitivity analysis, the LHS coverage is preferable because
it ensures the model is evaluated in every region of the parameter space.

\solnote{Part (c) is the quantitative check: students should compute the
actual mean and std and report how close they are to the theoretical values.
The LHS improvement in mean accuracy comes from the fact that the sample
mean of a stratified sample is an unbiased estimator with lower variance
than the sample mean of a random sample (stratification theorem).}

\end{solution}

%% ----------------------------------------------------------
\begin{solution}{9.3}{L3}

\solanswer

\solpart{a-b} \textbf{LHS + PRCC computation.}

Using $N = 200$, $\pm50\%$ around nominal values:
\begin{verbatim}
p = sir_nominal();
lb = 0.5*[p.k, p.beta, p.tau, p.L];
ub = 1.5*[p.k, p.beta, p.tau, p.L];
rng(42); X = lhs_sample(200, 4, lb, ub);
Y = arrayfun(@(i) sir_I_max(X(i,:), p), 1:200)';
[rho, pval] = compute_prcc(X, Y, {'k','beta','tau','L'});
\end{verbatim}

Expected PRCC values (approximate):
\begin{center}
\begin{tabular}{lcc}
\toprule
Parameter & PRCC & $p$-value \\
\midrule
$k$     & $+0.85$ & $< 0.001$ \\
$\beta$ & $+0.84$ & $< 0.001$ \\
$\tau$  & $+0.72$ & $< 0.001$ \\
$L$     & $-0.04$ & $> 0.5$  \\
\bottomrule
\end{tabular}
\end{center}

\solpart{c} \textbf{PRCC bar chart.} Sorted: $k \approx \beta > \tau \gg L$ (similar to SI ranking).

\solpart{d} \textbf{Strongest negative PRCC: $L$.}
$L$ (mean lifespan) has a negative PRCC: longer lifespans reduce peak infections.
In plain language: populations with longer lifespans have smaller background
mortality rates, which means the susceptible pool is replenished more slowly,
and a larger fraction of the population is already immune/recovered by the
time of peak infections. The effect is very small ($|\text{PRCC}| \approx 0.04$,
not statistically significant).

\solpart{e} \textbf{Comparison with local SI tornado plot.}
Both methods rank $k$ and $\beta$ highest, $\tau$ third, $L$ negligible.
The PRCC ranking agrees with the local SI ranking for this model.
Agreement is expected because peak infections $I_{\max}$ is monotone in
each parameter over the $\pm50\%$ uncertainty range (the model is
not near any threshold), satisfying the PRCC assumption.

\solnote{For the SIR model with $R_0 = 2.1$ and $\pm50\%$ variation,
the PRCC and local SI rankings should agree well. Discrepancies would
appear if: (a) $\beta$ variation includes $R_0 < 1$ (non-monotone at threshold),
or (b) parameter interactions are strong (nonlinear effects). Students
should note that agreement here is specific to this model and range.}

\end{solution}

%% ----------------------------------------------------------
\begin{solution}{9.4}{L4}

\solanswer

$Y = 2Q_1^2 + Q_2$, $Q_1, Q_2 \sim \mathcal{U}(0,1)$, independent.

\solpart{a} \textbf{Variance decomposition.}
$\E[Y] = 2\E[Q_1^2] + \E[Q_2] = 2/3 + 1/2 = 7/6$.
$\Var(Y) = 4\Var(Q_1^2) + \Var(Q_2)$.
$\Var(Q_1^2) = \E[Q_1^4] - (\E[Q_1^2])^2 = 1/5 - 1/9 = 4/45$.
$\Var(Q_2) = 1/12$.
$\Var(Y) = 4(4/45) + 1/12 = 16/45 + 1/12 = 64/180 + 15/180 = 79/180$.

$D_1 = \Var[\E(Y|Q_1)] = \Var[2Q_1^2 + 1/2] = 4\Var(Q_1^2) = 16/45$.
$D_2 = \Var[\E(Y|Q_2)] = \Var[2/3 + Q_2] = 1/12$.

\solpart{b} \textbf{Sobol indices.}
$S_1 = D_1/\Var(Y) = (16/45)/(79/180) = (16/45)(180/79) = 64/79 \approx 0.810$.
$S_2 = D_2/\Var(Y) = (1/12)/(79/180) = (180/12)/79 = 15/79 \approx 0.190$.
$S_1 + S_2 = 79/79 = 1.0$. Yes, they sum to 1, because the model is
\emph{additive} (no interaction terms), so the full variance is explained
by the first-order indices.

\solpart{c} \textbf{Total indices.}
For an additive model: $S_{T_1} = S_1 = 64/79$, $S_{T_2} = S_2 = 15/79$.
Total indices equal first-order indices when there are no interactions.

\solpart{d} \textbf{Adding interaction $Q_1 Q_2$: $Y' = 2Q_1^2 + Q_2 + Q_1 Q_2$.}
Without computing: adding $Q_1 Q_2$ introduces an interaction term.
$S_1$ decreases (some variance is now shared with $Q_2$); $S_2$ decreases
for the same reason. The total indices $S_{T_1} > S_1$ and $S_{T_2} > S_2$
because each parameter contributes to the interaction term.

\textit{Verification:}
$\Var(Y') = 4\Var(Q_1^2) + \Var(Q_2) + \Var(Q_1 Q_2) + 2\text{Cov}(\ldots)$
(with $\text{Cov}$ terms vanishing by independence).
$\Var(Q_1 Q_2) = \E[Q_1^2 Q_2^2] - (\E[Q_1 Q_2])^2 = (1/3)(1/3) - (1/2)^2(1/4)
= 1/9 - 1/16 = 7/144$.
$D_{12} = \Var(Y') - D_1' - D_2' - $ (remaining interaction variance).
Computing exactly: $S_1' < 0.810$, $S_2' < 0.190$, $S_1' + S_2' < 1$.
The interaction term $D_{12} > 0$ accounts for the remainder.

\solnote{Part (b)'s exact sum-to-1 for the additive model is a crucial
result: for any purely additive model, all first-order indices sum to 1
and total indices equal first-order indices. Part (d) demonstrates the
effect of interactions without requiring exact computation, testing
conceptual understanding over algebraic facility.}

\end{solution}

%% ----------------------------------------------------------
\begin{solution}{9.5}{L4}

\solanswer

Using \texttt{run\_global\_sa\_sir} from the repository with $N = 2000$.

\solpart{a} \textbf{Sobol indices (nominal parameters, $\pm50\%$ range).}

Approximate values from Jansen estimators:
\begin{center}
\begin{tabular}{lcccc}
\toprule
Param & $\hat{S}_i$ & $\hat{S}_{T_i}$ & $\hat{S}_{T_i} - \hat{S}_i$ \\
\midrule
$k$     & $0.38$ & $0.43$ & $0.05$ \\
$\beta$ & $0.38$ & $0.43$ & $0.05$ \\
$\tau$  & $0.20$ & $0.25$ & $0.05$ \\
$L$     & $0.00$ & $0.00$ & $0.00$ \\
\bottomrule
\end{tabular}
\end{center}

\solpart{b} \textbf{Parameters with $S_{T_i} \gg S_i$.}
For all three active parameters, $S_{T_i} - S_i \approx 0.05$, indicating
moderate interaction effects (primarily $k$-$\beta$ and $k$-$\tau$ interactions).
No parameter has $S_{T_i} \gg S_i$ for this model at this range;
interactions are present but modest. This is consistent with the SIR model
being a smooth, moderately nonlinear ODE.

\solpart{c} \textbf{Sums.}
$\sum S_i \approx 0.96 < 1$ (interaction terms absorb the remaining $\sim4\%$).
$\sum S_{T_i} \approx 1.11 > 1$ (interactions are double-counted in total indices).
For a purely additive model, both sums would equal 1. The observed values
confirm the presence of small but nonzero interactions.

\solpart{d} \textbf{Best parameter to measure precisely.}
To reduce $\Var(I_{\max})$ most, measure the parameter with the largest
$S_{T_i}$ (total index), since $S_{T_i}$ measures the variance that
would remain if all uncertainty in parameter $i$ were eliminated, including
interactions. Here $k$ and $\beta$ are tied at $S_{T_i} \approx 0.43$.
Measuring either $k$ or $\beta$ precisely reduces $\Var(I_{\max})$
by $\approx 43\%$ (more precisely, the variance conditional on knowing $k$ exactly
is $(1 - S_{T_k})\Var(I_{\max}) \approx 0.57\,\Var(I_{\max})$).

\solnote{Part (d) is a key policy application of total indices: the total
index, not the first-order index, answers ``which parameter measurement
would most reduce output uncertainty?'' First-order indices answer
``which parameter contributes most to variance independently?'' Both
are important, and students should be able to distinguish them.}

\end{solution}

%% ----------------------------------------------------------
\begin{solution}{9.6}{L4}

\solanswer

Using $r = 15$ Morris trajectories, $p = 4$ levels ($\Delta = 4/6 = 2/3$),
$\pm50\%$ parameter ranges.

\solpart{a} \textbf{Morris measures.} Approximate values from \texttt{morris\_screening}:
\begin{center}
\begin{tabular}{lccc}
\toprule
Param & $\mu^*$ & $\sigma$ & Classification \\
\midrule
$k$     & $\sim 80$ & $\sim 20$ & Important, nonlinear \\
$\beta$ & $\sim 78$ & $\sim 19$ & Important, nonlinear \\
$\tau$  & $\sim 42$ & $\sim 14$ & Important, nonlinear \\
$L$     & $\sim 1$  & $\sim 0.5$& Negligible \\
\bottomrule
\end{tabular}
\end{center}

\solpart{b} \textbf{Diagnostic plot.}
The $\mu^*$ vs.\ $\sigma$ plot shows $k$ and $\beta$ in the upper-right
(large $\mu^*$ and $\sigma$: important and nonlinear), $\tau$ in the
middle-right, and $L$ near the origin (negligible).

\solpart{c} \textbf{Most and least influential.}
Most: $k$ (or $\beta$, tied). Least: $L$.

\solpart{d} \textbf{Reduced Sobol (3 parameters, $L$ fixed).}

With $L$ fixed: Sobol cost $= N(m+2) = 2000\times5 = 10{,}000$ evaluations
vs.\ full $2000\times6 = 12{,}000$. Cost reduction: 17\%.
Approximate results: $S_k \approx 0.43$, $S_\beta \approx 0.43$, $S_\tau \approx 0.24$;
total indices similar. Fixing $L$ does not materially change the other
indices (since $L$ was negligible). The Morris screening correctly identified
that $L$ could be dropped.

\solnote{The Morris cost here is $15\times(4+1) = 75$ evaluations
vs.\ $12{,}000$ for full Sobol — a $160\times$ cost reduction for the
screening phase. This illustrates Morris's value as a cheap first step
that identifies which parameters can be dropped from the expensive Sobol analysis.}

\end{solution}

%% ----------------------------------------------------------
\begin{solution}{9.7}{L5}

\solanswer

\textbf{This exercise requires the student to choose and describe a specific
published model.} The model answer below uses the Chitnis dengue model~\cite{chitnis2008determining}
as an example; any 8+ parameter epidemic model with published parameter estimates is acceptable.

\solpart{a} \textbf{Model, parameters, uncertainty ranges.}
Chitnis et al.\ (2008) dengue model: 8 state variables, 14 parameters including
mosquito biting rate $b_v$, transmission probabilities $\beta_{vh}$ and $\beta_{hv}$,
mosquito and human mortality rates, and mean durations. Uncertainty ranges:
$\pm50\%$ for biological parameters (based on reported ranges in Table 2 of the paper);
$\pm30\%$ for demographic parameters (population data are more precisely known).

\solpart{b} \textbf{QOI justification.}
Primary: $R_0$ (determines whether dengue can persist in the population —
the most policy-relevant threshold). Secondary: equilibrium human infection
prevalence $\bar{I}_h/N_h$ (determines endemic burden for healthcare planning).

\solpart{c} \textbf{PRCC vs.\ Sobol choice.}
Use \textbf{Sobol indices}. The dengue model is moderately nonlinear and
has known interaction between biting rate $b_v$ and transmission probabilities
($R_0 \propto b_v^2\beta_{vh}\beta_{hv}$, so interactions are structural).
PRCC assumes monotone relationships; near threshold or at high $R_0$,
this assumption may not hold for all parameters. Sobol indices capture
interactions and do not require monotonicity.

\solpart{d} \textbf{Sample size $N$.}
Use $N = 2000$ for each Sobol matrix. Run the convergence check:
compute Sobol indices at $N = 200, 500, 1000, 2000$; report the
value of $N$ at which the first-order indices change by $< 2\%$ as $N$ doubles.
For a 14-parameter model, the total evaluation cost is $N(14+2) = 32{,}000$ runs.
If each run takes 2 minutes, this is $\approx 4.4$ days of computation on one core.

\solpart{e} \textbf{Presentation format.}
Two tornado plots of Sobol $S_{T_i}$ (one per QOI), with parameters
labeled as interventions: ``Reducing mosquito biting rate by $X\%$ through
bed nets reduces $R_0$ by approximately $Y\%$.'' A two-panel figure
showing which parameters appear in both tornado plots (co-optimization)
and which are QOI-specific. Use traffic-light color coding: red for
high-priority, green for low-priority intervention targets.

\solpart{f} \textbf{Potential misleading results.}
If the model's best-fit parameters give $R_0$ near 1 (endemic/non-endemic
threshold), Sobol indices at those nominal values will be dominated by
parameters near the threshold. A 2\% uncertainty in $\beta_{vh}$ might
shift the model from endemic to non-endemic, giving extremely large apparent
sensitivity for $\beta_{vh}$. Recommendation: compute Sobol indices at
three nominal parameter sets ($R_0 = 1.5$, $2.0$, $3.0$) and report
whether the parameter rankings are stable across this range.

\solnote{This is an open-ended design exercise. Full credit requires
all six components with operationally specific content. Vague statements
(``choose the right sample size,'' ``present clear figures'') without
specifics receive partial credit. The model choice affects all answers;
a reasonable choice with internally consistent answers earns full credit
even if different from the example above.}

\end{solution}

%% ----------------------------------------------------------
\begin{solution}{9.8}{L6}

\solanswer

\textbf{Methodological concerns with Sobol analysis ignoring correlation
among three parameters.}

\solpart{a} \textbf{Mathematical problem.}

Standard Sobol indices assume parameter \emph{independence}: the matrices
$A$ and $B$ are sampled independently from the joint distribution $p(\vec{p})
= \prod_j p_j(p_j)$. When three parameters are correlated, the correct
joint distribution is $p(\vec{p}) = p(p_1, p_2, p_3) \times \prod_{j>3} p_j(p_j)$,
which cannot be factored. Sampling independently from marginal distributions
and then compositing (the Saltelli method) generates non-physical parameter
combinations (e.g., high phosphorus loading but low nitrogen loading, which
does not occur in the real watershed system). The resulting Sobol indices
estimate sensitivity to artificially independent variations, not to the
actual coupled variations that the system experiences.

\solpart{b} \textbf{Effect on each correlated parameter's index.}

\textit{Phosphorus loading $P$:} The standard $S_1(P)$ likely \textit{overestimates}
the true first-order index. When $P$ and $N$ are positively correlated,
the actual variation of $P$ is accompanied by variation in $N$. The standard
estimator attributes only the $P$ variation to $S_1(P)$; the interaction
that would occur under the true distribution is ignored, artificially
inflating the apparent individual contribution of $P$.

\textit{Nitrogen loading $N$:} Same logic applies: $S_1(N)$ likely overestimates.

\textit{Water residence time $T_w$:} If $T_w$ is correlated with both $P$ and $N$
(longer residence = more loading), $S_1(T_w)$ may be \textit{underestimated}
because the positive correlation means that varying $T_w$ while holding $P$
and $N$ fixed understates the total effect of residence time in the real system.

\solpart{c} \textbf{Minimal correction.}

Use a copula-based sampling method that preserves the observed correlations.
The Iman--Conover rank correlation constraint method is the simplest:
(1) generate independent LHS samples; (2) apply the Iman--Conover algorithm
to induce the observed rank correlations among the three parameters;
(3) compute Sobol indices using the correlated sample matrices.
This adds one matrix transformation step to the existing workflow with
no additional model evaluations.

\solpart{d} \textbf{Why sum-to-1 is not a valid independence check.}

The Sobol first-order indices sum to 1 only for an additive model
with independent parameters. With correlated parameters, the indices
can still sum to approximately 1 even with incorrect independence assumptions,
because the sum reflects the total explained variance, which equals 1 by
construction (the variance is fully accounted for by some combination of
first-order and interaction terms). The sum being 1 says nothing about
whether the individual index attributions are correct — it is a property
of the decomposition algorithm, not of the physical model or the sampling design.

\solnote{This is a rigorous L6 exercise testing critical evaluation of
a real methodological error. Part (d) is the subtlest: students often
think ``indices sum to 1'' is a validity check, but it is not — it is
simply a consequence of the ANOVA decomposition. Full credit requires
a clear mathematical explanation, not just ``the check doesn't work.''}

\end{solution}
